% Laget automatisk av elevark/lag_elevark.py fra opplegg/repetisjon-relativitet.html. Endringer her blir overskrevet.
\documentclass[11pt]{article}
\usepackage{elevark}
\begin{document}
\elevarkhode{Fysikk 2 · Spesiell relativitetsteori}{Repetisjon: spesiell relativitetsteori}{https://fysikklab.no/opplegg/repetisjon-relativitet.html}{Opplegget på nett}

Den spesielle relativitetsteorien bygger på to enkle postulater, men følgene er overraskende. Her er en oversikt over begrepene og formlene, og blandede oppgaver med løsningsforslag.

\section*{Hvem måler hva?}

\begin{center}\small
\begin{tabularx}{\linewidth}{YYY}
\toprule
\textbf{Størrelse} & \textbf{Egenverdien måles} & \textbf{Sett fra et system der gjenstanden beveger seg} \\
\midrule
Tid & av en klokke som er med på begge hendelsene (egentid $t_{0}$) & lengre: $t = γt_{0}$ (tidsforlengelse) \\
Lengde & i systemet der gjenstanden er i ro (egenlengde $L_{0}$) & kortere i fartsretningen: $L = L_{0}/γ$ (lengdekontraksjon) \\
Samtidighet & to hendelser kan være samtidige i ett system & men ikke i et annet som beveger seg langs linja mellom dem \\
\bottomrule
\end{tabularx}
\end{center}

\textbf{Postulatene:} 1) Fysikkens lover er de samme i alle treghetssystemer. 2) Lysfarten i vakuum er den samme for alle observatører, uansett hvordan lyskilden beveger seg.

\section*{Formler}

\begin{itemize}
\item Lorentzfaktoren: $γ = 1/\sqrt{1 - v^{2}/c^{2}}$. $γ \geq 1$, og $γ$ går mot uendelig når $v$ går mot $c$.
\item Tidsforlengelse: $t = γt_{0}$. Lengdekontraksjon: $L = L_{0}/γ$.
\item Bevegelsesmengde: $p = γmv$. Hvileenergi: $E_{0} = mc^{2}$. Total energi: $E = γmc^{2}$. Kinetisk energi: $E_{\mathrm{k}} = (γ - 1)mc^{2}$.
\item Energi og bevegelsesmengde: $E^{2} = (pc)^{2} + (mc^{2})^{2}$.
\end{itemize}

Konstanter: $c = 3{,}00 \cdot 10^{8}\,\mathrm{m/s}$, $m_{\mathrm{e}}c^{2} = 0{,}511\,\mathrm{MeV}$, $m_{\mathrm{p}} = 1{,}673 \cdot 10^{-27}\,\mathrm{kg}$, 1~lysår er avstanden lyset går på ett år.

\section*{Oppgaver}

Oppgavene bruker fart som en brøkdel av $c$. Regn med minst tre gjeldende siffer i $γ$.

\begin{oppgaveliste}
\oppgave \fase{forklar}{Begrep}Er påstandene sanne eller usanne? Begrunn svaret.

a) En astronaut i et raskt romskip merker at klokka hennes går sakte.\newline b) Lys fra en lykt på et romskip som flyr med $0{,}5c$, har farten $1{,}5c$ sett fra jorda.\newline c) To hendelser som er samtidige for én observatør, er samtidige for alle.\newline d) Et romskip trekkes sammen i fartsretningen, men ikke på tvers av den.\newline e) Med en sterk nok motor kan et romskip nå lysfarten.
\svar
\oppgave \fase{test}{Regn}Regn ut $γ$ for $v = 0{,}60c$, $0{,}80c$ og $0{,}99c$. Hvor stor må farten være for at $γ$ skal bli 2? Sjekk med Lysklokka.\footnote{\url{fysikklab.no/sim/lysklokke.html}}
\svar
\oppgave \fase{test}{Regn}Myoner lages 15~km over bakken og har farten $0{,}995c$. Gjennomsnittlig levetid i ro er 2,2~µs. Hvor lang tid bruker de ned, målt fra jorda og målt av myonene? Hvor mange levetider tar turen i hvert system? Hvor tykk er atmosfæren sett fra myonene? Sjekk med simuleringen.\footnote{\url{fysikklab.no/sim/myoner.html}}
\svar
\oppgave \fase{test}{Regn}Et romskip reiser til en stjerne 4,37~lysår unna med farten $0{,}90c$. Hvor lang tid tar reisen målt fra jorda og målt om bord? Hvor lang er avstanden målt fra romskipet? Sjekk med simuleringen.\footnote{\url{fysikklab.no/sim/romreise.html}}
\svar
\oppgave \fase{test}{Regn}Et tog har egenlengden 200~m og kjører gjennom en tunnel med egenlengde 180~m med farten $0{,}60c$. Hvor langt er toget sett fra tunnelen? Er toget helt inne i tunnelen et øyeblikk? Hvor lang er tunnelen sett fra toget? Sjekk med simuleringen.\footnote{\url{fysikklab.no/sim/tog-og-tunnel.html}}
\svar
\oppgave \fase{test}{Regn}Et elektron akselereres fra ro gjennom en spenning på 1,0 MV. Finn den kinetiske energien, den totale energien, $γ$ og farten. Hvilken fart ville du fått uten relativitetsteori? Sjekk med simuleringen.\footnote{\url{fysikklab.no/sim/elektron.html}}
\svar
\oppgave \fase{test}{Regn}Finn bevegelsesmengden til et proton med farten $0{,}80c$. Hvor mange ganger større er den enn den klassiske verdien mv? Se Konstant kraft og lysfarten.\footnote{\url{fysikklab.no/sim/relativistisk-bevegelsesmengde.html}}
\svar
\oppgave \fase{utfordring}{Sammensatt}Sola stråler ut $3{,}83 \cdot 10^{26}\,\mathrm{W}$. Hvor mye minker massen til sola per sekund? Sola har massen $1{,}99 \cdot 10^{30}\,\mathrm{kg}$. Hvor stor del av massen mister den på 10 milliarder år?
\svar
\end{oppgaveliste}

\end{document}
